\zlabel{triptych:concise:commentary:start}
\section{Why Cosmic Upheaval Commands Joy: Judgment, the Marian Cloud, and Interior Grace}
\runninghead{Cosmic Upheaval, Cloud, and Grace}

Three textual difficulties stand in Luke~21:25--33. Right after men wither in fear, Jesus bids disciples lift up their heads (vv.~26--28), and speaks of steadfast heavenly powers (\latin{virtutes caelorum}, v.~26) being ``moved'' as the Son of Man comes ``in a cloud'' (\latin{in nube}, v.~27). Where Matthew~24:32 names only the fig tree, Luke adds ``and all the trees'' (\latin{et omnes arbores}, v.~29).

\subsection{Shaken Powers, the Cloud, and the Command to Rejoice}

St.~Gregory the Great (\work{Homiliae in Evangelia}~1.1--6) and St.~Ambrose of Milan (\work{Expositio Evangelii secundum Lucam}~X.37--45) approach these verses from opposite sides. Gregory reads the cosmic convulsions literally to break worldly attachment before Judgment: ``Our Lord and Saviour wisheth to find us ready at His second coming. Therefore He telleth us what will be the evils of the world as it groweth old, that He may wean our hearts from worldly affections'' (\S1; Bute~1908). For Gregory, the ``powers of heaven'' moved in verse~26 are the angelic choirs---angels, archangels, thrones, dominations, principalities, and powers---who will appear visibly to demand an exact reckoning of what our invisible Creator now bears patiently (\S2). In this Sunday's Mass, the \latin{Credo} confesses that same Judge returning in glory, while the Trinity Preface joins the assembly to the unceasing praise of angels, archangels, cherubim, and seraphim.

Facing that reckoning, Gregory explains why verse~28 commands joy (\S3; Oxford~1843, preserved in St.~Thomas Aquinas, \work{Catena Aurea in Lucam} on Luke~21:28). Because ``head'' often signifies the mind, \latin{levate capita} means \latin{exhilarate corda} (``gladden your hearts''). Quoting James~4:4, that friendship with the world is enmity with God, Gregory asks why a traveller exhausted by a hard road would wish the journey never to end. He presses the warning from a sudden whirlwind (\latin{nudius tertius}, \latin{subito turbine}) that had struck Rome two days earlier (\S\S5--6; Bl.~Ildefonso Schuster, I, p.~322, inadvertently remembers it as an earthquake). Comparing the aging world (\latin{senescentem mundum}) to a bent, panting old man, Gregory recalls how that blast uprooted ancient trees, wrecked houses, and threw down churches, catching men planning the next morning's business. Such calamities are heralds running before coming wrath (\latin{sequentis irae praecones}, \S6).

St.~Ambrose (\work{Exp.~in~Luc.}~X.37--42) reads the same verses symbolically in an Incarnational and ecclesial perspective. For him, the darkened sun is Christ obscured by unbelief, and the moon is the Church unable to reflect His light whenever carnal vices or love of mortal life interpose themselves in persecution. The falling stars are patriarchs and saints lost to view, while men wither (\latin{arescentibus}) when baptismal dew dries up: \latin{aret perfidia, fides rorat} (``unbelief parches, faith bedews,'' X.37--38). In contrast to Gregory's judicial awe, Ambrose asks how blessed angels can be ``moved'' (X.39--40). Answering that Christ comes both across the world and within each soul (\latin{in singulis}), he explains that as the Lord of hosts draws near, heavenly powers are moved by an increase of grace (\latin{augmentum gratiae}, Ps.~33:6; 2~Cor.~3:16). The third interpretation develops that interior coming.

On \latin{venientem in nube} (``coming in a cloud,'' v.~27, X.41--42), Ambrose asks how the Lord who ``hath set his tabernacle in the sun'' (Ps.~18:6) could enter in storm rain. Scripture's clouds carry spiritual dew to temper divine brightness: Israel's pillar of cloud, and Isaiah and Ezekiel revealing the Trinity's holiness through cherubim and seraphim (Isa.~60:8), echoed in the Trinity Preface. Above all, that cloud is the Virgin Mary: ``He came also in a light cloud, incarnate of the Virgin; for the prophet saw Him coming like a cloud from the East'' (Isa.~19:1). Isaiah calls her \emph{light}, Ambrose adds, because ``no earthly vices weighed her down.'' He confirms the identification from Luke~1:35: ``Look upon the cloud on which the Holy Spirit came down, and the power of the Most High overshadowed.''

\subsection{The Fig Tree's Fruit and the Renewal of Creation}

Gregory and Ambrose diverge directly on verses~29--30. Gregory (\S3) turns the parable inside out: \latin{fructus mundi ruina est}---``the fruit of the world is its ruin.'' The world grows only to fall, and its visible decay announces the eternal summer of God's kingdom. Ambrose (X.43--45) notices Luke's unique ``and all the trees'' (\latin{et omnes arbores}) beside Matthew~24:32 and reads the fruit positively. For him, the fig tree is Israel and ``all the trees'' are all Gentile nations blossoming with faith ``so that every tongue may confess to God, the Jewish people also confessing'' before the summer harvest of resurrection. Yet Ambrose keeps a warning through a double figure (\latin{duplicem figuram}): beside fruitful faith stands the barren fig tree covered only in leaves, the disguise of sinners and Antichrist (Gen.~3:7, ``the garments of exiles from paradise'').

Inside the first interpretation, two patristic arguments balance Gregory's severe maxim. St.~Cyril of Alexandria (\work{Commentary on Luke}, Sermon~CXXXIX, Payne Smith pp.~653--654; leaving vv.~29--33 unexpounded) marks verse~25 as the turn from Jerusalem's siege to the end of the age. He explains verse~28's ``redemption'' as bodily resurrection and cosmic refashioning: ``this earthly and infirm body shall put off corruption, and shall clothe itself with incorruption by Christ's gift,'' calling for watchful holiness (2~Pet.~3:10--11; Luke~21:36). On vv.~28--29 and~32--33, St.~Bede (\work{In Lucae Evangelium Expositio}~VI) notes that \latin{Amen} (``truly'') stays untranslated lest its solemnity grow cheap (\latin{ne vilesceret nudatum}), and takes \latin{generatio haec} (``this generation'') as either humanity or the Jewish people. On verse~33, where Gregory (\S4) contrasted sturdy sky and earth with passing speech to show the stability of Christ's Word, Bede distinguishes the lower airy heaven (\latin{aereum}) from the upper starry firmament (\latin{aethereum}) via 2~Pet.~3:5--7 and the Flood: judgment's fire sweeps only the lower atmosphere. Reconciling Luke~21:33 with Ecclesiastes~1:4 (``the earth standeth for ever'') via 1~Cor.~7:31, Apoc.~21:1, and Ps.~101:27, Bede teaches that heaven and earth pass away in outward form (\latin{per imaginem}), not in substance (\latin{per essentiam}). Cleansed by fire (\latin{per ignem tergetur}), creation is preserved in its own nature (\latin{in sua semper natura}), as winter yields to spring.

Bede applies to Luke~21:33 three arguments of St.~Augustine (\work{De civitate Dei}~XX), who does not expound it directly. In XX.14 (Apoc.~20:9--13, 1~Cor.~7:31--32), Augustine teaches transmutation, not destruction: ``The figure, therefore, passes away, not the nature.'' In XX.16 (Apoc.~21:1, 1~Cor.~7:31), he teaches universal fire renews the world to something better for renewed flesh. In XX.24, answering Porphyry from Ps.~101:26--28, he notes 1~Cor.~7:31, 1~John~2:17, and Matt.~24:35 (not Luke~21:33) speak more mildly than ``They shall perish,'' and argues via 2~Pet.~3:6--11 and the Flood that only the lower atmospheric heavens perish.

Gregory and Cyril confront literal public judgment in Luke~21:25--27, whereas Ambrose's Marian cloud and blossoming nations align with St.~Mary Major and link Christ's return to the Incarnation. Ambrose's \latin{in singulis} supplies the third interpretation's interior grace while remaining anchored in universal judgment.

\section{How Our Earth Yields Her Fruit: Virgin Soil, Ploughed Hearts, and the Eucharistic Pledge}
\runninghead{Our Earth Yielding Her Fruit}

The Mass draws its Sunday Alleluia (v.~8) and Communion (v.~13) from Psalm~84 (Heb.~85). Why ask God to \emph{show} His mercy, and how can human ``earth'' yield \emph{its own} fruit (\latin{fructum suum}) when all good comes from God? Here the second and third interpretations divide directly, partly over a Latin Psalter difference.

\subsection{The Holy Spirit Overshadowing Virgin Earth}

In the second interpretation, St.~Robert Bellarmine (\work{Explanatio} on Ps.~84:8) reads \latin{salutare tuum} directly of the Incarnate Son: ``Send us your Son for a Savior'' (Titus~2:11, 3:4). On verse~13, he compares virgin conception to wildflowers nourished solely by sun and rain from heaven. ``The Lord will give goodness,'' Bellarmine writes, means God ``will send his Holy Spirit from heaven, who will overshadow a virgin, and thus our land, which was never ploughed nor sown, and was altogether an untouched virgin, will yield her fruit''---Christ, ``the flower of the field, and the lily of the valleys'' (Cant.~2:1).

St.~Augustine grounds that Incarnational reading in verse~12 (\latin{Veritas de terra orta est}, ``Truth hath sprung out of the earth''). In \work{Enarrationes in Psalmos}~84.13 (Coxe~85.10), he asks: ``What is truth? The Son of God. What is the earth? Flesh,'' born of the Virgin Mary so righteousness might look down from heaven to justify sinners. In his Christmas Day sermon \work{Sermo}~185.1 (\latin{In Natali Domini}~II, PL~38, cols.~997--998), Augustine unfolds the same verse: ``Truth which is in the bosom of the Father sprang from the earth so as to be also in a mother's bosom \dots\ Truth for whom heaven does not suffice sprang from the earth so as to be laid in a manger.'' From that descent he issues a wake-up cry echoing Advent (Eph.~5:14): ``Awake, O man: for you God became man. Rise, you who sleep, and arise from the dead, and Christ shall enlighten you \dots\ You would have perished, had He not come.''

St.~Bernard of Clairvaux (\work{Missus est}, Hom.~1, Browne \& Nolan pp.~56, 58--60, esp.~p.~59) traces Psalm~84:13 across salvation history. Patriarchal promises are the seed, Old Testament figures and wonders are the blossom of Nazareth, and Christ Himself is the Fruit. He quotes Psalm~84:13 beside Psalm~131:11 (``Of the Fruit of thy womb I will set upon thy throne''). In his \work{Second Sermon for Advent} (p.~19), Bernard separately compares Mary's inviolate womb to an unploughed, unsown field yielding the divine Flower (Isa.~11:1; Cant.~2:1). At Communion, Sicard of Cremona (\work{Mitrale}~V.2, PL~213, col.~205) writes that by the goodness of the Holy Spirit ``our earth,'' the Blessed Virgin, conceives and yields God and man. Dom Prosper Gu\'eranger (\work{The Liturgical Year}~I, p.~137) likewise identifies ``our Earth'' with the Blessed Virgin ``made fruitful by the dew of heaven'' after Isaiah, adding that the divine Fruit has just given Himself by anticipation in Holy Communion.

\subsection{Prevenient Grace, Old Latin \textit{Suavitas}, and the Eucharistic Pledge}

In the same psalm commentary, St.~Augustine (\work{Enarr.~in~Ps.}~84.9, 14--16 [Coxe~85.6, 11--13]) takes a penitential path developed in the third interpretation. On verse~8 (\S9 [Coxe~\S6]), he explains that God must \emph{show} His mercy because fallen man boasts of his virtue. Once persuaded that every good is God's gift, a man ceases to be proud, is not lifted up, does not fall, stands firm, clings fast, abides, and enjoys God. Like a penniless wanderer told by a trustworthy messenger that he is a senator's son with a rich patrimony waiting at home, the believer asks in \latin{salutare tuum da nobis} to see Christ ``not as the Jews saw Him and crucified Him, but as the Angels see Him, and rejoice'' (1~John~3:2).

On verse~12 (\S14 [Coxe~\S11]), Augustine calls sinful man ``earth'' (Gen.~3:19): ``Confess thy sins, and truth shall spring out of thee,'' moving God's mercy to say, ``Let us spare this man, for he spareth not himself: let us pardon him, for he himself confesseth.'' At verse~13 (\S15 [Coxe~\S12]), Augustine's Old Latin Psalter reads \latin{suavitatem} (``sweetness'') where the Gallican Vulgate and 1962 Missal sing \latin{benignitatem} (``goodness''). Although his theology of grace accords with the Missal's text, his verbal argument rests on \latin{suavitas}. Bare commands cannot cure sin's allure until God pours interior sweetness in righteousness into hearts ``furrowed, as it were, with the plough of confession.'' Then ``thou who at first didst delight in drunkenness, shall rejoice in sobriety: [\dots]\ and thou who didst take delight in robbing, shalt delight now in giving: thou whom shows delighted, shalt delight in prayer.'' In verse~14 (\S16 [Coxe~\S13]), confession thus opens the Baptist's highway for Christ's footsteps (Matt.~3:3). St.~Leo the Great (\work{Sermon}~17.1; cf.~\work{Sermon}~12.4) preaches the same Advent husbandry: fasting and almsgiving cultivate ``the ground of our heart'' to yield ``the germ of righteousness and the fruit of love.''

Directly on this Mass's Communion, Bl.~Ildefonso Schuster (\work{The Sacramentary}~I, pp.~322--323) reads Psalm~84:13 as thanksgiving for the Eucharist, ``a pledge of his infinite goodness'' refreshing arid hearts with grace. William Durandus (\work{Rationale}~VI.3, fol.~261v) takes \latin{benignitatem} as good affection (\latin{affectum bonum}) and ``our earth'' (\latin{terra nostra}) as our flesh yielding good works. In the first interpretation's watchfulness, the same divine goodness sustains faithful endurance through the world's twilight. These expositions agree in doctrine while resting on distinct identifications of ``our earth'' and ``goodness'' or ``sweetness.''

\section{Waking from Sleep and Putting On Christ: Battle-Array, Incarnation, and Broken Habit}
\runninghead{Waking from Sleep and Putting On Christ}

Following St.~Paul's teaching that love fulfills the Law (Rom.~13:8--10), Romans~13:11--14a turns on three textual details. Paul declares that salvation is nearer now than when believers first believed (v.~11), and pairs ``works of darkness'' (\latin{opera tenebrarum}) with ``armour of light'' (\latin{arma lucis}, v.~12) rather than ``works of light.'' In addition, the 1962 Epistle stops at verse~14a (\latin{sed induimini Dominum Jesum Christum}), omitting verse~14b.

In the first interpretation, St.~John Chrysostom (\work{Homilies on Romans}~24) preaches the passage before Resurrection and Judgment. Beginners often slacken, so Chrysostom compares Christians to stadium runners pushing harder once they sight the finish and prize. To illustrate ``the night is passed, and the day is at hand,'' he pictures neighbors hearing the swallow before dawn and waking one another in the dark so sunrise will not catch them yawning in bed. Paul writes ``armour of light,'' Chrysostom explains, because steel armor is heavy and grim, whereas virtue guards the Christian while shining brighter than a sunbeam. Combat becomes a festival where the believer steps out ``at once soldier and bridegroom.'' Cornelius a Lapide (\work{Comm.~in~omnes~D.~Pauli~Epistolas}, pp.~195--198) treats \latin{arma lucis} as offensive and defensive virtues against the world, the flesh, and the devil (p.~196), and traces Paul's rejected carousing to the Greek \emph{komos}, nocturnal revelry (p.~197). In verse~14, Chrysostom follows an ascent from works to armor to putting on the King Himself, dwelling within as Way, Bridegroom, Food, High Priest, and Friend. Preaching on Luke~21:34 right after this Gospel, St.~Leo the Great (\work{Sermon}~19.1) calls for Advent fasting, prayer, and almsgiving to keep hearts from gluttony and worldly cares before that sudden Day.

The second and third interpretations read putting on Christ from two other angles. In the second, William Durandus (\work{Rationale}~VI.3, fol.~261r) draws the Incarnational exchange from verse~14a: ``for this reason the Son of God became man, that man might become a son of God.'' He also pairs (\work{Rationale}~VI.3, fol.~261v) the Epistle's ``our salvation is nearer'' (\latin{propior est nostra salus}) with the Gospel's ``the kingdom of God is at hand'' (\latin{prope est regnum Dei}). In the third, St.~Augustine (\work{Confessions}~VIII.12.29) recounts how Romans~13:13--14 broke his divided will in the Milan garden. Hearing \latin{tolle, lege} (``take up and read''), he read Paul's renunciation of revelry, impurity, and faction and felt ``a light, as it were, of security'' scatter every doubt. Yet his release turned on reading through verse~14b (``and make not provision for the flesh, to fulfil the lusts thereof''), whereas the 1962 Epistle omits that half-verse and closes directly on Christ's name. All three link this summons to the Collect's plea against inward peril (\latin{ab imminentibus peccatorum nostrorum periculis}). The first hears detachment before the Judge. The second hears the direct cry to the Son (\latin{Qui vivis}) echoing Psalm~79:3, where Bellarmine (\work{Explanatio} on Ps.~79:2--4) hears a plea to rouse God's hidden might. The third rests on Christ's protection and deliverance, as Schuster notes (I, p.~321).

\section{Lifting the Soul Along God's Ways and Paths: Briers, Humility, and the Ascent of Obedience}
\runninghead{God's Ways and Paths}

Psalm~24 (Heb.~25) recurs at the Introit (vv.~1--3, 4b), Gradual (vv.~3b, 4b), and Offertory (vv.~1--3). Across those chants, lifting the soul (\latin{Ad te levavi}) stands beside dread of shame (\latin{non erubescam}, ``let me not be ashamed''; \latin{non confundentur}, ``shall not be confounded'') and mocking foes. Twice before the Gospel, the assembly also asks for God's ``ways'' (\latin{vias}) and ``paths'' (\latin{semitas}), singing Gallican \latin{demonstra mihi} at the Introit and Old Latin \latin{notas fac mihi} at the Gradual.

In the first interpretation, St.~Robert Bellarmine (\work{Explanatio} on Ps.~24:1--3) hears a soul that has found only ``briers and thorns'' in creatures tearing itself from earthly ties to trust God alone against visible and demonic foes. It is steadied by the patriarchs' assurance that none who wait in sure hope are put to shame. On verses~4--5, Bellarmine adds that the psalmist asks to be taught ``not speculatively, but practically; that is to say, he asks for such grace as may move his will to observe the commandments cheerfully.'' Repeating \latin{Ad te levavi} at the Offertory entrusts that watchfulness to God at the altar. Rupert of Deutz (\work{De divinis officiis}~III.2, PL~170, col.~58) reads the repetition as the Church's continuing plea until she reaches the heavenly Jerusalem. Bl.~Ildefonso Schuster (I, p.~322) calls Psalm~24 ``the typical psalm of the first Sunday of Advent.'' In the second interpretation, John Beleth (\work{Summa}~LV, PL~202, col.~62) hears in the Introit the Church looking from afar toward the Nativity. At the stational basilica of St.~Mary Major, Schuster (I, p.~321) hears humanity ``cast down, yet full of hope'' asking to be led along the way of truth to Bethlehem.

The third interpretation examines the grammar of verses~1--4 through St.~Augustine and St.~Thomas Aquinas. Asking who speaks, Augustine (\work{Enarr.~in~Ps.}~24.1--4) answers with the whole Christ (\latin{totus Christus}): ``Christ speaks, but in the person of the Church,'' raising a soul trodden down by carnal desires. He traces shame back to Adam, who trusted himself to become like God and fell to ``fearing death from the smallest insect.'' Hidden spiritual foes still whisper flattery (``Well done, well done'') to lure believers off the narrow paths (\latin{semitae}) of Matthew~7:14.

St.~Thomas Aquinas (\work{Postilla super Psalmos}~24, Parma vol.~14, pp.~229--230) maps an ordered ascent. After St.~John Damascene, lifting the soul is prayer itself (\latin{ascensus intellectus in Deum}), joined to unhesitating trust (Jas.~1:6; Heb.~4:16). Shame (\latin{erubescentia}) springs from defeat, frustrated hope (Rom.~5:5), or past sin (Rom.~6:21). Commenting on Gallican \latin{qui sustinent te}, Aquinas explains patient endurance by the Roman chant's own verb: \latin{qui cum longanimitate eum expectat}, ``who waits for Him with longanimity'' (Hab.~2:3). On the omitted verse~4a, he distinguishes ruinous confusion, which is damnation, from salutary confusion, the shame of penance that brings grace (Ecclus.~4:25). On verse~4, Aquinas comments on Gallican \latin{Vias tuas, Domine, demonstra mihi, et semitas tuas edoce me} (``Show me thy ways, O Lord, and teach me thy paths''), not the Gradual's Old Latin \latin{notas fac mihi}. Dividing the active life into two grades, he defines \latin{viae} as the public highway of the commandments (\latin{praecepta}, centered in charity, John~14:15; Matt.~19:17), whereas a \latin{semita} is a narrower shortcut (\latin{quaedam abbreviatio communis viae}), the evangelical counsels (\latin{consilia}, Jer.~6:16; Isa.~40:3). Under \latin{Vias tuas} (``Thy ways''), because man does not know these ways by himself, the psalmist asks God to show them (\latin{ostende eas mihi}) both to the intellect and to the affections. Aquinas then turns under \latin{Et semitas} (``And Thy paths'') to the evangelical counsels, before treating verse~5's contemplative guidance (\latin{Dirige me in veritate tua, et doce me}, ``Direct me in thy truth, and teach me'').

\section{The Altar Between the Two Comings: Secret, Postcommunion, and the Threefold Advent}
\runninghead{The Altar and the Threefold Advent}

\subsection{St.~Bernard's Middle Coming and Bl.~Ildefonso Schuster's Altar Prayers}

St.~Bernard of Clairvaux (\work{Sermones in Adventu}~3.3--4, pp.~25--26; 5.1--3, pp.~40--42) names Christ's hidden visitation in grace the Middle Coming (\latin{adventus medius}). In Sermon~3.3--4 he distinguishes Christ's coming to men, His invisible coming into men (John~14:23), and His coming against men. In Sermon~5.1 he places that middle advent in spirit and virtue between the first coming in flesh and weakness (our Redemption) and the third coming in glory and majesty (our Life). He calls it our Repose and Consolation, ``the way, so to speak, by which we must travel from the first to the last.'' To keep the Word so this coming takes hold (5.2--3), Bernard compares it to bread. Left in a cupboard it grows stale, or thieves steal it, or mice nibble it, whereas bread eaten and digested passes into one's strength. So the Word must pass from mind into affections and will. In Sermon~4.1 (pp.~32--33), he urges believers to hold both visible comings together---``how much we have received by the first, how much we are promised at the second''---like the Bridegroom's two arms.

Bl.~Ildefonso Schuster (\work{The Sacramentary}~I, pp.~321--323) does not invoke Bernard's schema of a Middle Coming. On the Gospel (p.~322) he sets Bethlehem directly beside the second coming in glory, yet his exposition of the altar prayers describes the same present work of grace. In the Secret (\latin{Haec sacra nos}, p.~322), Schuster hears ``the ancient idea that runs through the Roman epiklesis,'' asking the sacrifice to purify our minds so we may worthily receive its Author (\latin{principium}). At Communion (pp.~322--323), he reads the Eucharist as refreshing arid hearts. In the Postcommunion (\latin{Suscipiamus}, p.~323), which he calls ``the true \latin{Eucharistia} in its etymological sense,'' the heavenly bread prepares worshippers for ``the approaching feast of our common ransom.'' William Durandus (\work{Rationale}~VI.3, fol.~259r) likewise notes this Postcommunion's preparation for the coming feast.

Across the three interpretations, these prayers and the governed Ordinary carry complementary accents. In the first, worshippers approach the returning Judge (\latin{principium}) in violet watchfulness beside the angelic praise of the Trinity Preface. Even while \latin{Gloria in excelsis} is silent, the Sunday Alleluia sounds because Advent is ``partly of rejoicing and partly of mourning'' (\latin{partim exultationis, et partim moeroris}, Durandus~VI.2, fol.~256v; compare Beleth, \work{Summa}~LXIV, PL~202, col.~71, \latin{partim laetitiae}, ``partly of gladness''). Rupert of Deutz (\work{De divinis officiis}~III.2, PL~170, col.~58) explains that the silent \latin{Gloria} signifies both the age that still awaited the joy of peace and the Church's present waiting for the consummation of salvation. In the second, the assembly approaches the divine Source of Incarnation and Eucharist at St.~Mary Major, praying for \latin{reparationis nostrae ventura solemnia} while withholding \latin{Gloria} so it may be sung more devoutly on Christmas night (Durandus~VI.2, fol.~256r). In the third, read beside Bernard's Middle Coming, the assembly receives present cleansing through the Secret, Communion, and Postcommunion before the 1962 Mass concludes with \latin{Ite, missa est} and the Last Gospel (John~1:1--14).

\subsection{Shared Ground, Real Disagreements, and the Four Senses}

\begin{comparisontable}{Point of Comparison}{Strict Judge}{Incarnational \& Marian}{Conversion \& Altar Bridge}
\textbf{Question} & How does cosmic ruin command joy, and what fruit does the fig tree yield? & Why station Advent~I at St.~Mary Major with a cloud, all trees, and fruitful earth? & How does grace heal sinners today so human soil bears fruit?\\
\textbf{Luke~21:25--33} & Literal heralds; angelic judges; ironic ``fruit of the world is its ruin'' (Gregory); resurrection (Cyril); fiery renewal \latin{per imaginem} (Bede, v.~33, applying Augustine, \work{De~civ.~Dei}~XX.14, 16, 24). & Symbolic darkening; \latin{nube} = prophets and Virgin Mary; ``all the trees'' = Israel and all nations blossoming (Ambrose). & Advent \latin{in singulis} increasing grace (Ambrose); Advent sobriety (Leo, Luke~21:34).\\
\textbf{Ps.~84:8 \& 13} & Cry for Savior in twilight; goodness sustaining endurance. & \latin{Salutare} = Son (Bellarmine); \latin{benignitas} = Spirit and untouched \latin{terra nostra} = Mary yielding Christ (Bellarmine, Sicard); Truth from Virgin flesh and Eph.~5:14 wake-up cry on v.~12 (Augustine, \work{Sermo}~185.1); Mary ``made fruitful by the dew of heaven'' (Gu\'eranger); Christ as Fruit (Bernard, \work{Missus est}, Hom.~1). & Mercy cures pride (Augustine, v.~8); Old Latin \latin{suavitatem} waters confessed hearts (Augustine, v.~13); Eucharistic pledge (Schuster); good works (Durandus).\\
\textbf{Ps.~24:1--4} & Tearing from earthly briers (Bellarmine) into unashamed watchfulness. & Church looking toward the Nativity (Beleth); humanity led to Bethlehem (Schuster). & \latin{Totus Christus} humbled from pride (Augustine); ascent and commandments (\latin{viae}) vs.\ counsels (\latin{semitae}, Aquinas).\\
\textbf{Rom.~13:11--14a} & Runners nearing the prize; glittering \latin{arma lucis}, ``soldier and bridegroom'' (Chrysostom); \emph{komos} (Lapide). & Son became man so man might become son of God; ``salvation is nearer'' paired with ``kingdom is at hand'' (Durandus). & Decisive break with carnal habit (Augustine, \work{Confessions}~VIII.12.29; Missal stops at v.~14a).\\
\textbf{Spiritual Senses} & \textit{Alleg.:} ruin as fruit, angelic judges. \textit{Moral:} detachment, sobriety. \textit{Anag.:} Parousia, bodily incorruption, earth renewed \latin{per imaginem}. & \textit{Alleg.:} Marian cloud, virgin earth, blossoming nations. \textit{Moral:} guarding baptismal dew (\latin{fides rorat}). \textit{Anag.:} summer harvest of nations. & \textit{Alleg.:} \latin{totus Christus}, Middle Coming, purifying Mysteries. \textit{Moral:} confession, humility, precepts to counsels. \textit{Anag.:} pure approach to \latin{principium}.\\
\end{comparisontable}

All three interpretations share the Collect's confession of sin's peril, St.~Paul's call to wake and put on Christ, trust in Psalm~24, longing for mercy in Psalm~84, Luke~21's coming Son of Man, and the Trinity Preface. Two exegetical divergences separate them without creating doctrinal conflict. On Luke~21:29--30, St.~Gregory reads the fig tree's fruit ironically as the world's decay, whereas St.~Ambrose reads Luke's ``and all the trees'' as Israel and the nations blossoming in faith. On Psalm~84:13, St.~Robert Bellarmine, Sicard of Cremona, and Dom Prosper Gu\'eranger identify ``our earth'' with the Virgin Mary yielding Christ, alongside St.~Bernard naming Christ the Fruit (\work{Missus est}, Hom.~1). By contrast, St.~Augustine (working from Old Latin \latin{suavitatem}) expounds the verse of sinners' hearts ploughed by confession, while Bl.~Ildefonso Schuster applies it to Holy Communion refreshing arid souls. Taken together, the three interpretations correspond to Christ's threefold coming (\latin{triplex adventus}), summarized by John Beleth (\work{Summa}~LXIV, PL~202, col.~72) and preached by St.~Bernard (\work{In~Adventu}~3.3--4, 5.1). Christ came once in human flesh through the Virgin Mary, comes daily within human souls by the Holy Spirit and grace, and will come openly in majesty to judge and renew creation. In the third interpretation, that daily coming in grace is set beside Bl.~Ildefonso Schuster's Eucharistic exposition of the Secret, Communion, and Postcommunion.
